% fixture: amsart/amsart-sections
% tier: T3
% engine: document
% covers: amsart's \section (centered small capitals, the number in \mdseries with a period and \enspace), run-in \subsection (bold) and \subsubsection (italic) with \@addpunct (after an uppercase letter's period too), an empty \subsection title (a bold number, then \@secnumfont \bfseries for the later numbers too; \textup's \/), a heading right after a heading, starred forms; amsart's ~ (\unskip\nobreak\ \ignorespaces)
% issue: #360
\documentclass{amsart}
\begin{document}
\section{Introduction}
Typesetting is the composition of text by means of arranging physical type or its digital equivalents. Stored letters and other symbols are retrieved and ordered according to the orthography of a language.

\subsection{Background}
The breaking of paragraphs into lines is a small part of this craft, but it decides how evenly the words are spread.

\subsection{Why it matters?}
And how easily the eye moves from one line to the next.

\subsubsection{A detail}
More text follows here in the subsubsection to fill a line or two of the page with words.

\subsection{Results in the U.S.}
An uppercase letter before the period: one more is added.

\section{Next}
\subsection{Right after}
A subsection right after a section, no penalty or skip between them.

\subsection{}
An empty subsection title: the number in bold and a period.

\section*{Unnumbered}
Text after an unnumbered section, as in Section ~2 and in Lemma~ 3: amsart's ties take the spaces around them.
\section{A much longer section title that will need to break over two lines in the centered heading}
Final words.
\end{document}
